\documentclass[10pt,a4paper]{amsart}
\usepackage{amsmath,amssymb,amsthm,graphicx}
% Retained minimal preamble for the Stacks Project chapter more-algebra.tex (More on Algebra).
% The chapter file pinned-original-source.tex carries only \input{preamble} on its first line, so no
% preamble text lives inside the chapter itself.  The Stacks Project uses one single shared preamble for
% every chapter; the snapshot stages it identically as sources/stacks-chow/preamble.tex and
% sources/stacks-etale-cohomology/preamble.tex.  The block below is copied mechanically, line for line,
% from sources/stacks-chow/preamble.tex lines 14-18 (the Xy-pic setup) and 154-217 (the native theorem
% declarations together with the mathematical macros).  Nothing was added, removed or redefined; the
% wrapper class and the amsmath/amssymb/amsthm/graphicx base packages are supplied by the assembler.
% Chapter: more-algebra.tex -- More on Algebra.
% For commutative diagrams we use Xy-pic
\usepackage[all]{xy}

% We use 2cell for 2-commutative diagrams.
\xyoption{2cell}

% Theorem environments.
%
\theoremstyle{plain}
\newtheorem{theorem}[subsection]{Theorem}
\newtheorem{proposition}[subsection]{Proposition}
\newtheorem{lemma}[subsection]{Lemma}

\theoremstyle{definition}
\newtheorem{definition}[subsection]{Definition}
\newtheorem{example}[subsection]{Example}
\newtheorem{exercise}[subsection]{Exercise}
\newtheorem{situation}[subsection]{Situation}

\theoremstyle{remark}
\newtheorem{remark}[subsection]{Remark}
\newtheorem{remarks}[subsection]{Remarks}

\numberwithin{equation}{subsection}

% Macros
%
\def\lim{\mathop{\mathrm{lim}}\nolimits}
\def\colim{\mathop{\mathrm{colim}}\nolimits}
\def\Spec{\mathop{\mathrm{Spec}}}
\def\Hom{\mathop{\mathrm{Hom}}\nolimits}
\def\Ext{\mathop{\mathrm{Ext}}\nolimits}
\def\SheafHom{\mathop{\mathcal{H}\!\mathit{om}}\nolimits}
\def\SheafExt{\mathop{\mathcal{E}\!\mathit{xt}}\nolimits}
\def\Sch{\mathit{Sch}}
\def\Mor{\mathop{\mathrm{Mor}}\nolimits}
\def\Ob{\mathop{\mathrm{Ob}}\nolimits}
\def\Sh{\mathop{\mathit{Sh}}\nolimits}
\def\NL{\mathop{N\!L}\nolimits}
\def\CH{\mathop{\mathrm{CH}}\nolimits}
\def\proetale{{pro\text{-}\acute{e}tale}}
\def\etale{{\acute{e}tale}}
\def\QCoh{\mathit{QCoh}}
\def\Ker{\mathop{\mathrm{Ker}}}
\def\Im{\mathop{\mathrm{Im}}}
\def\Coker{\mathop{\mathrm{Coker}}}
\def\Coim{\mathop{\mathrm{Coim}}}

% Boxtimes
%
\DeclareMathSymbol{\boxtimes}{\mathbin}{AMSa}{"02}

%
% Macros for moduli stacks/spaces
%
\def\QCohstack{\mathcal{QC}\!\mathit{oh}}
\def\Cohstack{\mathcal{C}\!\mathit{oh}}
\def\Spacesstack{\mathcal{S}\!\mathit{paces}}
\def\Quotfunctor{\mathrm{Quot}}
\def\Hilbfunctor{\mathrm{Hilb}}
\def\Curvesstack{\mathcal{C}\!\mathit{urves}}
\def\Polarizedstack{\mathcal{P}\!\mathit{olarized}}
\def\Complexesstack{\mathcal{C}\!\mathit{omplexes}}
% \Pic is the operator that assigns to X its picard group, usage \Pic(X)
% \Picardstack_{X/B} denotes the Picard stack of X over B
% \Picardfunctor_{X/B} denotes the Picard functor of X over B
\def\Pic{\mathop{\mathrm{Pic}}\nolimits}
\def\Picardstack{\mathcal{P}\!\mathit{ic}}
\def\Picardfunctor{\mathrm{Pic}}
\def\Deformationcategory{\mathcal{D}\!\mathit{ef}}

% --- appended 2026-10-12 (Chief): the project own definitions of the reference/slogan/history
% display environments, copied mechanically from the shared Stacks preamble lines 8-12
% (\usepackage{verbatim} plus the three \newenvironment lines, which the chapter preamble keeps
% verbatim from the project preamble). Nothing added, removed or redefined: these environments
% carry no mathematical content and the project itself typesets them as comments. Without these
% three lines a counted statement whose own text uses them fails with Environment slogan undefined.
\usepackage{verbatim}
\newenvironment{reference}{\comment}{\endcomment}
\newenvironment{slogan}{\comment}{\endcomment}
\newenvironment{history}{\comment}{\endcomment}
\usepackage{cleveref}
\crefname{originallemma}{original lemma}{original lemmas}
\Crefname{originallemma}{Original lemma}{Original lemmas}
\makeatletter
\newlabel{lemma-galois-inertia}{{09EE}{source}{}{sourcepointer-1}{}}
\newlabel{lemma-galois-inertia@cref}{{[originallemma][0][]09EE}{[1][0][]source}{}{}{}}
\makeatother
\crefname{originallemma}{original lemma}{original lemmas}
\Crefname{originallemma}{Original lemma}{Original lemmas}
\makeatletter
\newlabel{lemma-galois-galois}{{09ED}{source}{}{sourcepointer-2}{}}
\newlabel{lemma-galois-galois@cref}{{[originallemma][0][]09ED}{[1][0][]source}{}{}{}}
\makeatother
\title{The inertia character is equivariant for the action of the decomposition group}
\author{The Stacks Project Authors}
\date{stacks-project:more-algebra.tex:local-snapshot-20261012; GNU FDL 1.2 or later}
\begin{document}
\maketitle
\setcounter{subsection}{1}
\noindent\textit{Editorial note.} All delivered mathematical text is the author's own, copied from the pinned Stacks Project chapter file more-algebra.tex; no statement or proof step was rewritten, reordered or removed. The counted statement is the labeled lemma environment at source lines 33430-33439 (canonical tag 0BU5) and the counted proof is the proof environment at source lines 33441-33468, which immediately follows it and is its own proof.\par\medskip
\noindent\textit{Editorial note.} The statement and proof use the notation of the surrounding section for the decomposition group, the inertia group, the uniformizer and the residue field; those notations are introduced by the two lemmas the proof refers to, which are carried as navigation stubs.\par\medskip
\noindent\textit{Editorial note.} The preamble span runs to line 90, the current end of the retained preamble, which includes the appended environment definitions; the counted text does not use them, and the span is taken whole so that the file version is unambiguous.\par\medskip
\noindent\textit{Editorial note.} The source URL is the source-level Stacks Project URL, as required when more than one card is produced from one source; the per-claim tag is carried by claim\_key and by this note. The counted statement and proof cite no bibliography entry, so this card needs no bibliography.\par\medskip
\noindent\textit{Allowed standard prerequisites (editorial).}\begin{itemize}
\item The inertia group is normal in the decomposition group, and the decomposition group acts on the residue field by the map to its automorphism group. used at proof lines 33442-33446. Explains why the formula of the statement is meaningful and what the action on the right hand side is.
\item An inertia element acts on a uniformizer by multiplication by a unit, whose image in the residue field is the value of the inertia character. used at proof lines 33447-33453. Recalls the construction of the character, which the computation uses.
\item The relation between the unit attached to an element of the decomposition group and the unit attached to its inverse, together with the triviality of inertia elements modulo the uniformizer. used at proof lines 33454-33468. Supplies the substitutions that turn the value at the conjugate into the conjugate of the value at the original element.
\end{itemize}\medskip
\ifdefined\hypertarget\hypertarget{sourcepointer-1}{}\fi
\noindent\textit{Source reference.} Original lemma 09EE: sources/stacks-more-algebra/pinned-original-source.tex, source lines 33251--33251. The reference is the construction of the inertia character recalled at proof lines 33447-33453; the computation of the conjugate that proves the statement is inside the counted span.\par\medskip
\ifdefined\hypertarget\hypertarget{sourcepointer-2}{}\fi
\noindent\textit{Source reference.} Original lemma 09ED: sources/stacks-more-algebra/pinned-original-source.tex, source lines 33224--33224. The reference supplies the action of the decomposition group on the residue field invoked at proof lines 33444-33446; it is a pointer needed to make the formula meaningful and the rest of the argument is inside the counted span.\par\medskip
\begin{lemma}
\label{lemma-inertia-character}
With assumptions and notation as in Lemma \ref{lemma-galois-inertia}.
The inertia character $\theta : I \to \mu_e(\kappa(\mathfrak m))$
satisfies the following property
$$
\theta(\tau \sigma \tau^{-1}) = \tau(\theta(\sigma))
$$
for $\tau \in D$ and $\sigma \in I$.
\end{lemma}

\begin{proof}
The formula makes sense as $I$ is a normal subgroup of $D$
and as $\tau$ acts on $\kappa(\mathfrak m)$ via the map
$D \to \text{Aut}(\kappa(\mathfrak m))$ discussed in
Lemma \ref{lemma-galois-galois} for example.
Recall the construction of $\theta$. Choose
a uniformizer $\pi$ of $B_\mathfrak m$ and for
$\sigma \in I$ write $\sigma(\pi) = \theta_\sigma \pi$.
Then $\theta(\sigma)$ is the image $\overline{\theta}_\sigma$
of $\theta_\sigma$ in the residue field. For any $\tau \in D$
we can write $\tau(\pi) = \theta_\tau \pi$ for some unit $\theta_\tau$.
Then $\theta_{\tau^{-1}} = \tau^{-1}(\theta_\tau^{-1})$.
We compute
\begin{align*}
\theta_{\tau \sigma \tau^{-1}}
& =
\tau(\sigma(\tau^{-1}(\pi)))/\pi \\
& =
\tau(\sigma(\tau^{-1}(\theta_\tau^{-1}) \pi))/\pi \\
& =
\tau(\sigma(\tau^{-1}(\theta_\tau^{-1})) \theta_\sigma \pi)/\pi \\
& =
\tau(\sigma(\tau^{-1}(\theta_\tau^{-1}))) \tau(\theta_\sigma) \theta_\tau
\end{align*}
However, since $\sigma$ acts trivially modulo $\pi$ we see that
the product $\tau(\sigma(\tau^{-1}(\theta_\tau^{-1}))) \theta_\tau$
maps to $1$ in the residue field. This proves the lemma.
\end{proof}

\end{document}
